% 本文件由 LaTeX 排版 Agent 根据有效 translation.json 自动生成。
% author=Tertullian; work_id=0402; title=论妇女的服饰

\chapter{论妇女的服饰}

% structure: p0001 source=h1 target=omitted reason=page-title-already-rendered-by-parent

% structure: p0002 source=h2 target=section reason=relative-heading-rank; base=h2

\section{第一卷}

% structure: p0003 source=h3 target=subsection reason=relative-heading-rank; base=h2

\subsection{第一章 导言：妇女服饰的端庄，纪念罪因妇女进入世界}

如果在地上存在着如同在天上所盼望的那样大的信德，至亲爱的姊妹们，就没有一位，从起初\Quote{认识主}并学到关于自身（即女性）境况的真理之后，会渴望过于喜乐（且不说过于炫耀）的衣着样式；反倒会穿着谦卑的服装，宁愿显得卑微，像厄娃那样哀恸悔改地行走，为能藉着每一种补赎的装束更完全地补偿她从厄娃所继承的——即第一宗罪的耻辱，以及作为人类沦亡之因的憎恶。\Quote{你怀孕生子必受痛苦和焦虑；你要依恋你的丈夫，他要管辖你。} 你岂不知道你（每一个）就是厄娃？天主对你这性别的判决活在这时代：罪责也必然活着。\Emph{你}是魔鬼的门径；\Emph{你}是那（禁）树的开启者；\Emph{你}是天主法律的第一个背弃者；\Emph{你}是说服那魔鬼不敢攻击之人的那位。\Emph{你}如此轻易地摧毁了天主的肖像——人。因\Emph{你}的功绩——即死亡——连天主子也必须死去。而你还想在你那些皮衣之外装饰自己？来吧；如果从世界之初米利都人就剪羊毛，塞里斯人纺树木，提洛人染色，弗里吉亚人刺绣，巴比伦人织布，珍珠闪烁，缟玛瑙发光；如果金子本身也已带着贪欲从地里出来；如果镜子也早已如此大肆说谎，那么被逐出乐园的厄娃（已死的厄娃），我想也会贪图\Emph{这些}东西！\Emph{如今}，她（若想再活）就更不应渴求或知晓她\Emph{活着}时既未拥有也不认识的东西。因此，这些东西都是女人在定罪和死亡状态下的行装，仿佛为了增添她葬礼的排场而设立的。\par

% structure: p0005 source=h3 target=subsection reason=relative-heading-rank; base=h2

\subsection{第二章 女性装饰的起源，追溯至堕落的天使}

因为那些创始者自己也被定罪，归于死亡的刑罚——就是那些从天上冲向人间女子们的天使；于是这耻辱也加在女人身上。因为当他们向一个远比我们无知的时代揭示了一些隐藏很好的物质材料，以及若干未被充分揭示的科学技艺——如果说他们揭示了冶金术的操作，传播了草本的天然特性，颁布了巫术的力量，并探索了每一种好奇的技艺，甚至到了解释星辰的地步——他们就恰当地、仿佛特地为女人提供了那种女性炫耀的工具：宝石的光辉装饰项链，黄金的环圈压缩手臂，染料药物染色羊毛，以及那使眼睑和睫毛突出的黑色粉末。这些事物的性质可以暂且从它们教师的性质与状况来宣告：罪人绝不能展示或提供任何有助于正直的事物，非法恋人绝不能提供有助于贞洁的事物，叛离的灵体绝不能提供有助于敬畏天主的事物。如果这些东西被称为\Emph{教导}，邪恶的老师必然教导邪恶；如果作为\Emph{情欲的报酬}，就没有任何卑贱的事物其报酬是光荣的。但为何展示和赐予这些东西如此重要？难道女人没有光彩的物质原因和巧妙的恩宠技巧，就不能取悦那些尚未装饰、粗野、可以说原始粗糙的\Emph{男人}，却曾打动\Emph{天使}的心吗？还是说那些恋人若没有赠予被引诱与之结姻的女子任何（补偿性）礼物，就会显得卑鄙和——因无偿使用——无礼？但这些问题无法估量。拥有天使（为丈夫）的女人不可能再渴望更多；她们确实缔结了伟大的婚姻！当然，她们有时会想起自己从哪里堕落，并在情欲的炽热冲动之后仰望天空，于是用那自然美（作为邪恶之因）的卓越本身来回报，以使她们的幸运毫无益处；反而她们与那些天使一同，从纯朴和真诚转离，成为天主所憎恶的。他们确知一切炫耀、野心和藉肉欲取悦的爱好是\Emph{不}悦乐天主的。这些就是我们将要审判的天使；这些就是我们在圣洗中弃绝的天使；这些当然就是他们之所以配受人类审判的理由。那么，他们的\Emph{东西}与他们的\Emph{审判者}有什么关系？那些将要定罪的人与那些将被定罪的人有什么交往？我想，正如同基督与贝里雅耳不相干一样。我们有何种一致坐上那（未来）审判席，对那些人宣判，却又（如今）追求他们的礼物？因为你们（作为女人）也有同样应许的天使本性作为赏报，同男性一样的性别；主应许你们同样提升到审判的尊严。除非我们开始在此\Emph{预先}审判，预先定罪他们的\Emph{东西}——而这些我们将来要在他们\Emph{自身}上定罪——那么\Emph{他们}将审判并定罪\Emph{我们}。\par

% structure: p0007 source=h3 target=subsection reason=relative-heading-rank; base=h2

\subsection{第三章 论\WorkTitle{哈诺客预言}的真实性}

我意识到，\WorkTitle{哈诺客书}将这种（行动的）秩序归于天使，但有些人不接受它，因为它也不被列入犹太正典。我想他们认为，既然它在洪水之前发表，不可能安全地幸存于那场毁灭万物的普世灾祸。如果那是（拒绝它的）理由，让他们回想一下，洪水的幸存者\Person{诺厄}就是\Person{哈诺客}本人的曾孙；他当然从家喻户晓的声望和继承传统中听说过并记得关于他曾祖父的\Quote{在天主眼中蒙恩宠}以及他所有的宣讲；因为\Person{哈诺客}没有给\Person{默突舍拉}别的吩咐，只让他将这些知识传给后代。因此\Person{诺厄}无疑可能成功继任了（他）宣讲的托管；或者，如果情况并非如此，他也不会对天主（他的保护者）所安排的事物以及他自己家族的独特荣耀保持沉默。\par

如果（\Person{诺厄}）没有通过如此短的途径拥有这种（保守能力），那么（仍然）会有这种（考虑）来支持我们关于这部圣经（真实性）的主张：他同样可以在\Emph{圣神}的默感下，在它被洪水的暴力摧毁后，\Emph{重新编写}它，正如在\Place{耶路撒冷}被巴比伦人攻陷后，犹太文学的每一份文献通常被认为是通过\Person{厄斯德拉}恢复的。\par

但是，既然\Person{哈诺客}在同一部圣经中也宣讲了关于主的事，那么凡与我们有关的一切，我们绝对不可拒绝；我们读到：\Quote{凡适合建立德行的圣经，都是天主所默感的。}现在\Emph{犹太人}似乎正是因为（那个）理由而拒绝了它，正如几乎所有其他讲述基督的部分一样。当然，他们不接受一些谈论祂的圣经这一事实并不奇怪，因为他们甚至连亲自在他们面前发言的祂本人也不接受。除了这些考虑，还有\Person{哈诺客}在宗徒\Person{犹达}的书信中得到见证这一事实。\par

% structure: p0011 source=h3 target=subsection reason=relative-heading-rank; base=h2

\subsection{第四章 暂且不讨论作者问题，\Person{戴尔都良}提议根据事物本身的优劣来考虑。}

现在姑且认为妇女的虚荣没有被（其）创造者的命运烙上预先定罪的印记；除了他们背叛天堂和（他们的）肉性婚姻之外，不要归咎于那些天使：让我们考察事物本身的特性，以便我们也能觉察它们被热切追求的目的。\par

女性服饰包含两个观念——衣着和装饰。所谓\Quote{衣着}，我们指她们所谓的\Quote{女性雅饰}；所谓\Quote{装饰}，我们指适合称为\Quote{女性\Emph{丑}饰}的东西。前者被认为（包括）金银、宝石和衣物；后者包括对头发、皮肤以及那些吸引眼目的身体部位的照料。针对前者我们指控为野心，针对后者则指控为卖淫；因此，即使从（我们讨论的）早期阶段，你就能看到，在这些之中，哪一样适合天主的婢女——即\Emph{你的}操行，因为你基于不同的原则（与其他女人）被评估——也就是谦卑和贞洁的原则。\par

% structure: p0014 source=h3 target=subsection reason=relative-heading-rank; base=h2

\subsection{第五章 金银在起源和用途上并不优于其他金属。}

金银，作为世俗辉煌的主要物质原因，必然与它们所从出的东西（在本质上）相同：即土；（而土本身）显然比它们更光荣，因为只有当土在被诅咒的矿场的致命工场中，经过惩罚性的劳动而泪流满面地加工，并在那里将它的\Quote{土}这个名字留在火中之后，它才作为矿场的逃犯，从折磨变为装饰，从惩罚变为美化，从耻辱变为荣耀。但是铁、铜以及其他最卑贱的物质材料，享受着（与金银）等同的条件，无论是在地上起源还是冶金操作上；为的是在自然的评价中，金银的材料不会被判断为比它们高级一丝一毫。但如果金银是从\Emph{效用}的性质中获得其荣耀，那么，铁和铜胜过它们；它们的用途被（造物主）如此安排，以至于它们不仅自身行使更多、对人类事务更必要的功能，而且凭借自身的力量，在为更公正的事业服务中，也同样服务于金银的用途。因为不仅戒指是用铁做的，而且古人的记忆还保存着某些用铜制成的饮食器皿。让金银疯狂的富足自己去面对，如果它甚至被用来制造污秽用途的器具。无论如何，田地不是用金耕种的，船只也不是用银的强度连接的。没有锄头将金刃插入土中；没有钉子将银尖钉入木板。我不提整个生活的需要都依赖铁和铜这一事实；而那些富贵的材料本身，既需要从矿中挖出，而且在每一次（被投入的）使用中都需要锻造过程，没有铁和铜的辛勤力量就毫无用处。因此，我们已经必须判断，为什么如此高的尊荣归于金银，因为它们优先于那些不仅在起源上是它们的表亲，而且在效用上更为强大的物质材料。\par

% structure: p0016 source=h3 target=subsection reason=relative-heading-rank; base=h2

\subsection{第六章 论宝石和珍珠。}

然而，其次，那些与金子争高而傲慢的宝石，除了是些小卵石、石头和同一泥土的卑微微粒外，我还能将它们解释为何物？但它们对于奠基、筑墙、支撑门楣或加固屋顶都非必要。它们唯一懂得建造的，就是女人这种愚蠢的骄傲：因为它们需要缓慢的摩擦才能发光，巧妙的衬垫才能显美，精心的钻孔才能悬挂；它们与金子互相配合，以行娼妓般的诱惑。但无论野心从不列颠海或印度海中捞起何物，都是一种贝类，其\Emph{滋味}并不比——我不说牡蛎和海螺——甚至比巨蚌更令人愉悦。容我补充说，我知道那些贝类是海洋的甜美果实。但如果那种（来自远方的）贝类内部生了某种脓疱，那应当被视为它的缺陷而非荣耀；尽管它被称为\Quote{珍珠}，但必须理解为鱼身上某种坚硬、圆形的赘生物，而非其他。有些人还说，宝石是从\Emph{龙}的额上采摘的，正如鱼脑中有某种石质物。这竟然也是基督徒妇女所缺少的，好让她从蛇身上增添风韵！难道她就这样将脚踩在魔鬼的头上，同时却将取自他头上的装饰堆在自己的颈项上，甚至堆在自己的头上吗？\par

% structure: p0018 source=h3 target=subsection reason=relative-heading-rank; base=h2

\subsection{第七章. 唯独稀有使此类物品珍贵}

所有这些物品的优雅完全来自它们的稀有和异域性；简言之，在其本土范围内，它们并不被如此看重。丰裕总是自轻自贱的。有些野蛮民族，因为黄金本地盛产，习惯用金链锁住罪犯狱中的囚犯，让恶人满载财富——罪越重，财富越多。终于，竟然找到了一种方法，连金子都不被爱慕！我们在罗马也曾见过，当帕提亚人和玛代人，以及他们自己的同胞轻蔑地使用宝石时，我们的贵妇面前的宝石因受轻视而黯然失色——只是（\Emph{他们的}宝石）通常不用于炫耀。翡翠藏在他们的腰带里；佩在胸前的剑是唯一见证装饰剑柄的圆柱形宝石的；靴子上巨大的单颗珍珠甘愿从泥泞中拾起！总之，他们携带的饰物中，没有哪件比那件（本\Emph{不}该饰以宝石，若不让它显眼，或显眼只是为了表明它也被忽视的东西）更奢华的。\par

% structure: p0020 source=h3 target=subsection reason=relative-heading-rank; base=h2

\subsection{第八章. 同样的规则也适用于颜色。天主的受造物通常不得使用，除非用于祂所指定的目的}

同样，连那些野蛮人的仆人（因穿着类似的衣物）也使我们衣袍的色彩黯然失色；甚至他们的墙壁也轻蔑地使用推罗紫、紫罗兰色和庄严的皇家挂毯来替代绘画，而你们却费力地拆解它们并加以改造。在他们那里，紫色比红赭石更不值钱；（这是合理的，）因为衣服从非法的颜色中能获得什么正当的尊荣呢？祂自己没有产生的东西，不会令天主喜悦，除非祂\Emph{不能}命令羊生出紫色和天蓝色的羊毛！如果祂\Emph{能够}，那么显然祂\Emph{不愿}：天主不愿的事，当然不该被制造。因此，那些按\Emph{本性}不是从天主——本性的\Emph{创造主}——而来之物，并非至善。由此可见，它们来自\Emph{魔鬼}，来自本性的\Emph{败坏者}：因为倘若不是天主的，便没有其他来源；凡不是天主的，必定属于祂的敌手。但除魔鬼及其天使外，没有别的天主敌手。再者，如果\Emph{物质材料}来自天主，并不能立即断定人间\Emph{享受}它们的方式也是如此。要探究的不仅在于贝类从何而来，还在于它们被赋予了何种装饰领域，以及它们在何处展示其美。因为所有那些世俗娱乐的亵渎乐趣——我们已出版了一卷关于它们的论著——（甚至）连偶像崇拜本身，也都从天主的受造物中获取物质原因。然而，基督徒不应沉溺于赛车的疯狂、竞技场的暴行或舞台的丑恶，仅仅因为天主赐给人马、豹子和说话的能力；正如基督徒也不能安然行偶像崇拜，因为乳香、酒、喂养（祭火）的火和用作牺牲的动物都是天主的化工；甚至被崇拜的物体本身也是天主的（受造物）。因此，对于它们的积极使用，物质材料的\Emph{起源}源于天主，却\Emph{开脱}了那种（使用）与天主相违，犹如犯了世俗荣耀的罪过！\par

% structure: p0022 source=h3 target=subsection reason=relative-heading-rank; base=h2

\subsection{第九章. 天主的分配必须规整我们的欲望，否则我们成为野心及其伴随的邪恶的猎物}

因为，正如天主分配在某些个别土地上的特定事物，以及某一片特定的海域，彼此互不相干，它们要么被忽视，要么被渴望：被外邦人渴望，因为是稀罕之物；如果在自己的同胞中，则（合理地）被忽视，因为在\Emph{它们}那里，对于一种在本土是冷淡的光荣，没有如此炽热的渴望。然而，由天主随其意愿所安排的财产分配而产生的稀罕与异域风情，常常在陌生人眼中得宠，仅仅因为\Emph{没有}天主使别处固有的东西，就激发了\Emph{拥有}它的贪欲。由此又产生另一种恶习——\Emph{无节制}的拥有；因为，尽管也许\Emph{拥有}是允许的，但必须遵守一个限度。这（第二种恶习）将是野心；因此，它的名称也可以这样解释，因为它是由心中\Emph{环绕}的贪欲生出的，旨在追求光荣的欲望——一个伟大的欲望，当然，它既非由本性也非由真理所推荐，而是由一种邪恶的心灵激情——（即）贪欲。还有与野心和光荣相关的其他恶习。因此，它们也提高了物品的\Emph{代价}，以便借此为自己添加燃料；因为贪欲变得相应地更大，因为它对自己急切渴望的东西赋予了更高的价值。从最小的首饰盒里产出一份丰厚的祖产。一根丝线上悬挂着一百万塞斯特斯。一个纤细的脖颈环绕着森林和岛屿。薄薄的耳垂耗尽一笔财富；左手，每根手指，玩弄着几个钱袋。这就是野心的力量——（足以）在一个小小的身体上，而且还是女人的身体上，承载如此丰富财富的产物。\par

% structure: p0024 source=h2 target=section reason=relative-heading-rank; base=h2

\section{第二卷}

% structure: p0025 source=h3 target=subsection reason=relative-heading-rank; base=h2

\subsection{第一章 引言。端庄不仅要在其本质上，也要在其附属品上被遵守}

永生天主的婢女，我的同仆和姊妹，我同你们共享的权利——我，在那同仆和弟兄关系中是最卑微的——使我胆敢向你们致辞，当然不是出于情感，而是为你们得救的缘故为情感铺路。那救恩——不仅是女人的救恩，同样也是男人的——首先在于展现端庄。因为，由于圣神被接入我们内并使之成为我们的，我们都是\BibleQuote{天主的殿宇}，端庄是那殿宇的圣器管理员和女司祭，她不容许任何不洁或亵渎之物被引入（殿内），免得居住其中的天主被冒犯，完全离开那被玷污的居所。但当前场合，我们（要讲）的不是关于端庄本身，因为四面八方压迫我们的神圣诫命足以命令和要求端庄；而是关于与端庄相关的事，即你们应当如何行事。因为大多数女人（我信赖天主准许我，当然是为了我个人的指责，来指责一切），要么出于单纯的无知，要么出于掩饰，竟胆敢这样行事，仿佛端庄只在于肉体的（纯粹）完整和远离（实际的）淫乱，而不需要任何外在的东西——我是指服饰和装饰的安排，形态和光彩的刻意优雅：——她们在举止中穿戴着与外邦女子完全相同的外表，而外邦女子缺乏\Emph{真正}端庄的意识，因为在那些不认识天主——真理的守护者和主人——的人中，没有\Emph{任何}真实的东西。因为如果在外邦人中可以相信有任何端庄存在，那么显然它必然是不完全和无约束的，以至于尽管它在心里一定程度上积极保持自己，却仍允许自己松懈为衣着的放纵奢侈；正符合外邦人的乖戾，渴望那他们小心避免其后果的东西。简而言之，有多少人不热切渴望甚至看上去令陌生人愉悦？谁不因此小心把自己涂抹打扮，却否认自己曾是（肉欲）欲望的对象？然而，即使这也是外邦端庄的常见做法——即不是实际\Emph{犯}罪，但仍\Emph{愿意}犯罪；或者甚至不\Emph{愿意}，却仍然不完全拒绝——这有什么奇怪呢？因为一切不属于天主的都是乖戾的。因此，让那些女人留意吧，她们因为没‌有持守\Emph{全部}善，很容易甚至把她们所持守的与恶混杂。你们必须远离她们，如同在其他一切事上，也如此在你们的举止上；因为你们应当\BibleQuote{完美，如同你们的天父是完美的一样。}\par

% structure: p0027 source=h3 target=subsection reason=relative-heading-rank; base=h2

\subsection{第二章 完全的端庄将从一切引向罪的事物中戒绝，如同从罪本身戒绝。信赖与妄断的区别。如果我们自己安全，不可将诱惑置于他人面前。我们必须爱邻人如同自己}

你当知道，在完美的端庄——即基督信仰的端庄——眼中，来自他人的对自己的（肉身）欲望不仅不应被渴望，更当被你憎厌：第一，因为谋求将个人的容貌（我们知道它本是淫欲的引诱者）作为取悦手段，并非出自健全的良心：为何要在自己身上激起那邪恶的（情欲）？为何要吸引你所自认的陌路者？第二，因为我们不应为诱惑开路，那些诱惑因其紧迫，有时竟成就一种（恶行），而天主从祂的人身上驱逐这恶行；（或）至少藉着（设置）绊脚石，使心神极度骚乱。我们确实应当如此圣洁地行走，并带着如此完整的信仰实质，以致对自己的良心充满自信与安稳——\Emph{渴望}那（恩赐）能存留于我们内心直到终结，却不可\Emph{自恃}（必能如此）。因为自恃的人少忧虑；少忧虑的人少谨慎；少谨慎的人冒更多风险。恐惧是救恩的根基；自恃是恐惧的障碍。因此，担忧自己可能失足，比自信不会失足更为有益；因为担忧将引我们恐惧，恐惧引向谨慎，谨慎引向救恩。反之，若我们自恃，便既无恐惧也无谨慎来拯救我们。行事稳妥却不警醒的人，没有稳妥而坚固的安全；而警醒的人才能真正稳妥。愿主以祂的仁慈眷顾祂的仆人，使\Emph{他们}甚至得以\Emph{自恃}祂的美善！但我们为何成为邻人的（危险）？为何将私欲输入邻人？这私欲，若天主在\Quote{增补律法}中没有将它从实际的奸淫行为之中分离出来加以刑罚，我不知道祂是否会容许那使他人丧亡的人逍遥法外。因为那人一旦对你的美貌起了贪欲，并且内心已犯了私欲所指的（行为），他便丧亡了；而你成了毁灭他的利剑：因此，纵然你免于（实际）罪行，却免不了（随之而来的）憎恶；正如有人地产被盗，那（实际）罪行固然不会归咎于地主，然而庄园却蒙受耻辱，（而）地主本人也染上恶名。难道我们要打扮自己好让邻人丧亡吗？那么，诫命\Quote{你当爱你的邻人如你自己}又在哪里？\Quote{不只顾自己的事，也顾及邻人的事}？圣神的每一句话都不应只局限于当下谈论的主题，而应适用于每一用得上的场合。既然我们的自身利益和他人利益都牵扯到对最危险的（外在）美貌的追求，那么你该知道：不仅虚构和精心打扮的盛饰必须被你摒弃，就连天生丽质也当以隐藏和疏忽来抹杀，因为它同样对（旁观者的）眼目有危险。因为，尽管美貌不可\Emph{谴责}——它既是身体的一种福气，又是神圣塑造艺术的额外付出，可算是灵魂的一件华美外衣——却仍当\Emph{畏惧}，只因追求者的伤害和强暴：这种（伤害和强暴）甚至连信仰之父亚巴郎也因自己妻子的美貌而大为畏惧；依撒格也靠谎称黎贝加为妹妹而用侮辱换取了安全！\par

% structure: p0029 source=h3 target=subsection reason=relative-heading-rank; base=h2

\subsection{第三章 姑且承认美貌不足畏惧：但因其无用且虚荣，仍当回避}

现在姑且承认，优美的形貌无须畏惧，既不给所有者带来麻烦，也不使渴慕者毁灭，更不给同伴带来危险；可以认为它不暴露于诱惑，也不被绊脚石围绕：仅凭它对天主的众天使并非必要，就已足够。因为哪里有了端庄，那里美貌便无所作为；因为美貌的正当用途和果实本是纵欲，除非有人以为肉身的魅力还能收获其他。那些女人认为，在向\Emph{邻人}提供美貌所需之物时，她们也在为\Emph{自己}增添那（天生）的美貌，并在未得美貌时竭力追求，难道不是吗？有人会说：\Quote{那么，既然纵欲被排除、贞洁得以进入，我们岂不可单享受美貌的称赞，并在肉身的善中夸耀么？}愿那以\Quote{在肉身内夸耀}为乐的人自去审视。首先，按天主的诫命，我们这些\Emph{谦逊}的宣认者并不追求\Quote{光荣}，因为\Quote{光荣}是\Emph{高举}的实质，而\Emph{高举}与宣认\Emph{谦逊}不相称。其次，如果\Emph{一切}\Quote{光荣}都是\Quote{虚幻}和无知的，那么\Emph{在肉身内}的光荣，尤其对我们而言，岂不更是如此？因为即使\Quote{夸耀}是被允许的，我们也应希望自己的悦人之道在于圣神的恩宠，而非在于肉身；因为我们乃是属神之事的\Quote{追求者}。我们在哪些事上劳苦，也当在哪些事上喜乐。从我们期待得救的泉源，让我们采摘我们的\Quote{光荣}。显然，一个基督徒\Emph{也}会\Emph{在肉身内}\Quote{夸耀}，但那是当它为基督的缘故受撕裂时，为使灵魂在其中得冠冕，而非为使年轻人的眼睛和叹息追随它。因此，无论从哪方面看，这事在你身上都是多余的：若你没有，尽可鄙弃；若你有，尽可忽略。一个圣洁的女子，如果天生美貌，就不该给人如此大的（情欲）机会。确实，即使她如此美丽，也不应炫耀，反而应当掩盖它。\par

% structure: p0031 source=h3 target=subsection reason=relative-heading-rank; base=h2

\subsection{第四章　關於\Quote{討丈夫喜悅}的辯解}

好像我是在對外邦人說話，用外邦人的誡命向你們講話（一條共同的誡命），（我會說：）\Quote{你們只應當討你們丈夫的喜悅。}但是，你們越是不在乎討別人喜悅，就越能討他們的喜悅。蒙福的（姊妹們），不要憂慮：沒有一位妻子在自己丈夫眼中是\Quote{醜陋的}。當她被（他）選中（為妻）時，她已經\Quote{討}了他足夠的\Quote{喜悅}；無論是因容貌或品行而被稱讚。你們中誰也不必以為，如果她不在意外表，就會招致丈夫的憎恨和厭惡。每個丈夫所要求的都是貞潔；而信（主的丈夫）並不要求美貌，因為我們不被外邦人所認為的那些魅力所吸引；反之，不信（主的丈夫）甚至會因外邦人對我們那臭名昭著的看法而猜疑。那麼，你是為誰而妝飾你的美貌呢？若為信者，他並不要求；若為不信者，除非是樸素的美，他不會相信。你何必急於討好那猜疑的，或那並不希冀的呢？\par

% structure: p0033 source=h3 target=subsection reason=relative-heading-rank; base=h2

\subsection{第五章　某些服飾與儀容的修飾是許可的，某些是不許可的。顏料屬於後者}

當然，這些建議並不是要你們發展成完全粗野和狂放的外表；我們也不是要說服你們相信污穢和邋遢的益處；而是要說明個人修飾的限度、標準和適當尺度。不可越過那條界限，即簡單而適當的修飾所限定的欲望界限——那條界限是取悅天主的。因為那些用藥膏擦皮膚、用胭脂塗臉頰、用銻粉使眼睛突出的人，是得罪祂。我想，天主的塑造技藝令他們不悅！我想，他們在自己身上定罪、指責了萬物的創造者！因為他們修改、添加（祂的作品）時，就是在指責；這些添加當然是從對頭的工匠那裡取來的。那個對頭的工匠就是魔鬼。除了那因邪惡而改變人精神者，誰能指示改變身體的方法呢？無疑是他，他設計了這類巧妙的伎倆；這樣，在你們身上就可以顯明，你們在某種意義上是在對天主施暴。凡是天生的，都是天主的工程。凡是貼上去的（東西），都是魔鬼的工程。將撒殫的巧計疊加在天主的工程上，這是何等罪惡！我們的僕人不從我們個人的敵人那裡借任何東西；士兵不渴望從自己將帥的敵人那裡得到任何東西；因為從你所屬的那位對手那裡索取任何東西供自己使用，都是違犯。基督徒能從那惡者那裡得到任何幫助嗎？（如果那樣，）我不知道\Quote{基督徒}這名號是否還會繼續屬於他；因為他將屬於他所渴望從其教訓中學習的那一位。然而，這些事與你們的教訓和職業是何等背道而馳！戴一張假面具，與基督徒的名號何等不相稱，（你們）被命令要在各方面保持純樸！——在外表上撒謊，（你們）連用舌頭撒謊都不被許可！——追求別人的東西，（你們）被交付了戒除別人東西的誡命！——在臉色上犯姦淫，（你們）以謙遜為學問！思想吧，蒙福的（姊妹們），如果你們不在自己身上保持天主的形像，你們怎能遵守天主的誡命呢？\par

% structure: p0035 source=h3 target=subsection reason=relative-heading-rank; base=h2

\subsection{第六章　關於染髮}

我看見有些（婦人）用番紅花改變（她們）頭髮的（顏色）。她們甚至以自己的民族為恥，（恥）於她們的出身沒有把她們分配到日耳曼和高盧；於是，她們就把自己的（頭髮）轉移（到那裡）！用她們火焰色的頭為自己預測凶兆，惡啊，實在是極惡！她們以為那是優雅，實則是在玷污！不僅如此，化妝品的力量還會把毀壞燒進頭髮；即使不間斷地使用任何未加藥的水分，也會給頭部積存傷害；而太陽的溫暖，雖然對頭髮生長和乾燥如此有益，卻是有害的。什麼\Quote{優雅}能與\Quote{傷害}並存？什麼\Quote{美麗}能與\Quote{不潔}並存？一個基督徒婦人會把番紅花堆在頭上，如同在祭台上嗎？因為，凡慣常被焚燒以尊崇不潔之靈的，都似乎是祭獻——除非用於正當、必要和有益的用途，即天主的造物被賦予的用途。然而，天主說：\Quote{你們中間誰能將一根白頭髮變黑，或將黑頭髮變白？}於是她們反駁了主！她們說：\Quote{看哪，我們不去做白或黑，而是做成黃色——更有吸引力的優雅。}而那些後悔活到老年的人，甚至試圖把它從白色改成黑色！何等冒失！我們所渴望和祈求的年齡，卻在為自己臉紅！一場盜竊發生了！我們曾在其間犯罪的那青年歲月，被嘆息懷念！清醒的機會被糟蹋了！願如此大的愚昧遠離智慧之女！老邁越是試圖隱藏自己，就越會被發現。看，在你們頭髮的青春中，有一個真正的永恆！這裡我們有一種\Quote{不朽壞}可以\Quote{穿上}，以進入神聖王國所應許的新居！你們直奔向主，你們急於擺脫這個最邪惡的世界，對你們而言，接近（自己的）終點是難看的！\par

% structure: p0037 source=h3 target=subsection reason=relative-heading-rank; base=h2

\subsection{第七章　關於以其他方式過分裝飾頭髮及其對救恩的影響}

再者，花在\Emph{整理}头发上的所有劳苦对救恩有何益处？为什么不让你的头发得享安宁，时而束起，时而松开，时而培养，时而稀疏？有些人急于把头发卷成卷发，有些人则让它松散飘垂；并非出于良好的纯朴：此外，你们还附上一些我不知是什么奇形怪状的精巧编织的假发；时而像生皮头盔的样子，仿佛护头的鞘和冠盖；时而向后（拉）成堆在颈上。奇怪的是，居然没有（公开）违反主的命令！经上已宣布，没有人能增加自己的身量。然而，你们却给自己的\Emph{重量}增添某种卷轴或盾形饰物，堆在脖子上！如果你们对这巨大之物不感到羞耻，至少要为玷污而感到羞耻；恐怕你们是在把一个圣洁和基督徒的头部，配上别人头上可能不洁、可能有罪、注定下地狱的蜕皮。不，倒要彻底从你们\Quote{自由}的头上赶走这一切装饰的奴役。你们徒劳地努力显得装饰华丽；徒劳地请来所有最精巧的假发制造者。天主要求你们\Quote{蒙头}。我相信（祂这样做）是怕有些人的头被看见！哦，但愿在基督徒欢欣的\Quote{那日子}，我这最可怜的人（尽管在你们脚后跟以下）也能抬起头来！我（那时）将看见你们是否带着铅粉、胭脂和番红花，以及所有那些头饰的炫耀而升起：是否就是这样打扮的妇女被天使提到空中迎接基督？如果这些（装饰物）\Emph{现在}是好的，并且属于天主，那么它们\Emph{那时}也会出现在复活的身上，并认出各自的部位。但除了纯粹的肉身和神魂，没有什么能复活。因此，凡不是以神魂和肉身的形式复活的，都被定罪，因为它不出于天主。你们要从被定罪的事上戒绝，即使在今日。在今天，愿天主看见你们正如祂\Emph{那时}将看见你们一样。\par

% structure: p0039 source=h3 target=subsection reason=relative-heading-rank; base=h2

\subsection{第8章 关于个人修饰的评论不排除男性}

当然，现在我作为一个男人，好像是嫉妒妇女，要把她们完全赶出自己的（领域）。在我们男人方面，是否也有一些事情，鉴于我们因敬畏天主而应保持的节制，是不被允许的呢？如果这是真的（事实如此），那么在男人中，为了妇女的缘故（正如在妇女中为了男人的缘故），由于本性的缺陷，植入了取悦的意愿；而且，我们这个性别承认自己特有的欺骗性的形式技巧——（例如）过分尖利地修剪胡须；这里那里拔掉胡须；剃光（嘴）的周围；整理头发，用染料掩饰灰白；剃去全身初生的柔毛；用某种女性化的化妆品把（每根毛发）固定在原位；用某种粗糙的粉末或其他东西来平滑身体其余部分；然后，进一步利用每个机会照镜子；焦虑地凝视镜子——然而，一旦天主的认识结束了所有通过感官诱惑取悦的意愿，所有这些都被视为轻浮、敌视端庄而被摒弃。因为哪里有天主，哪里就有端庄；那里就有节制作为她的助手和盟友。那么，没有她的工具手段，即没有节制，我们怎能实践端庄呢？再者，我们怎能将节制应用于履行端庄的（职能），除非在外貌、面容以及整个人的总体印象中表现出严肃来标志我们的举止？\par

% structure: p0041 source=h3 target=subsection reason=relative-heading-rank; base=h2

\subsection{第9章 应避免着装和个人的过度修饰。引自格前7的论据。}

因此，关于服饰，以及所有其余你们自我修饰的累赘之物，你们也必须同样地修剪和削减过于华丽的铺张。因为在\Emph{脸上}表现出节制、质朴和与神圣训导完全相称的简朴，却用奢靡可笑的排场和精致来装点身体的\Emph{其他}部分，这有什么益处呢？这些排场与纵欲之事的关系何等密切，它们又如何干扰端庄，这一点很容易从以下事实看出：正是借助服饰的联合帮助，它们才卖弄了个人容貌的恩宠；所以显然，如果没有这些排场，它们便会使那恩宠变得无用和无感，仿佛被解除武装、支离破碎。另一方面，如果天生丽质不足，外在装饰的辅助帮助便能提供一种恩宠，仿佛是出于它自身固有的力量。事实上，那些最终得享宁静并退入端庄之港湾的人生时光，服饰的光辉与尊严会将它们引离（那宁静和那港湾），并藉由食欲的诱惑\Emph{扰乱}严肃，那诱惑借着衣着的挑逗魅力来弥补年岁的寒冷。首先，有福的（姊妹们），（当心）你们不要使用娼妓和卖淫式的装束和衣物；其次，如果你们中有人因财富、出身或过去的尊贵的需要，被迫如此华丽地出现在公众面前，以致看起来似乎尚未获得智慧，则要小心调和这种邪恶；免得在需要的借口下，无节制地放纵自由的放任。因为如果你们不克制对财富和雅致的享受，而这些享受又如此倾向于\Quote{光荣吗？}，你们怎能履行我们（学派）所教导的谦逊呢？因为光荣的习性向来是\Emph{高举}，而不是\Emph{谦卑}。\Quote{为什么，我们不可以使用自己的东西吗？}谁禁止你们使用呢？然而，必须依照宗徒，他告诫我们\BibleQuote{要使用这个世界，就像没有滥用它一样；因为这世界的形态正在消逝。}并且\BibleQuote{购买的人要如此行动，仿佛他们并不拥有。}为什么这样？因为他已立下前提，说：\BibleQuote{时间已经缩短了。}那么，如果由于时代的紧迫，他明确地表明甚至连妻子本身也要如同\Emph{没有}一般地拥有，那么他对她们这些虚浮的装饰又会持什么看法呢？此外，难道不是有许多人确实\Emph{这样做}，并为天国的缘故自阉，自发地放弃如此光荣且（如我们所知）被允许的欢愉吗？难道不是有些人禁止自己使用\BibleQuote{天主的受造之物,}戒绝酒和肉食，这些享受并不接近任何危险或忧虑；但他们甚至在节制的饮食中也将灵魂的谦逊祭献给天主吗？因此，在（领受）救恩训导的知识之前，你们也充分使用了你们的财富和精致；充分削减了你们嫁妆的果实。我们是\BibleQuote{万世的终结已来到，结束了它们的行程。}的人。我们是在世界之前由天主预定的，在时间的极终（出现）。因此，天主训练我们，为的是惩戒，而且（可以说是）阉割这个世界。我们是在一切事上受割损的——属灵和属肉体的；因为我们在灵性上和肉体上都割损世俗的原则。\par

% structure: p0043 source=h3 target=subsection reason=relative-heading-rank; base=h2

\subsection{第十章。\Person{戴尔都良}再次论及所有这些装饰和美化之物的起源问题。}

毫无疑问，是天主指示了用草汁和贝类体液染羊毛的方法！他在命令宇宙生成时，竟忘了为（生产）紫色和朱红色的羊发出命令！也是天主精心设计了那些本身轻而薄、但价格沉重的衣物的制造工艺；天主制造了如此华贵的金饰具用于束发或分开发丝；天主引入了耳朵上精巧切割的伤口，如此珍视对自己作品的折磨和纯洁婴儿的酷刑，使他们在最初的呼吸中就学会受苦，以便从那些为钢刃而生的身体疤痕上悬挂我不知道什么（珍贵的）颗粒，这些颗粒，正如我们可以清楚看到的，帕提亚人将它们插入鞋子以代替饰钉！然而即使是金子本身，其\Quote{光荣}使你们着迷，却也服务于某个民族（正如外邦文献告诉我们的）用作锁链！因此，构成这些事物之良好品质的并非其内在价值，而是其稀有性，这一点是如此真实；此外，通过犯罪的天使所引入的人工技巧来加工它们的过度劳动，这些天使同时还是这些物质本身的启示者，与稀有性相结合，激发了它们的珍贵，从而在女性方面产生了拥有那种珍贵的贪欲。但是，如果同样的天使，既揭示了这类物质及其魅力——即金子，和璀璨的宝石——并教人如何制作它们，随后又教导他们，在其他（训示）中，涉及眼影粉和羊毛染色的（功效），却已被天主定罪，正如\Person{哈诺客}告诉我们的，那么当我们因那些天使的\Emph{事物}而喜悦时，我们怎能中悦天主呢？那些天使正因这些事激起了天主的\Emph{愤怒}和\Emph{报复}？\par

如今，姑且承认天主预见了这些事，也允许了它们；依撒意亚并未指责任何紫色衣袍，也未禁止任何卷发，更未谴责任何新月形颈饰；但我们不可像外邦人那样自欺，以为天主仅仅是创造者，而非同时是祂所造之物的俯察者。因为，如果我们冒昧地推测，这一切确实起初由天主预备并安置在世界中，为的是现在可以考验祂仆人的纪律，使得\Quote{使用}的许可成为进行\Quote{节制}实验的手段，那将是何等更有益、更谨慎的做法！明智的家长难道不是故意向仆人提供并允许某些东西，以试探他们是否以及如何使用这些被允许之物——是否诚实、有节制？然而，完全禁戒、甚至对主人的宽容怀有健康畏惧的仆人，岂不更值得称赞？因此，宗徒也说：\Quote{凡事都可行，}他说，\Quote{但不全有益。}一个人若对\Emph{合法}之事怀有敬畏，岂不更容易畏惧\Emph{不}合法之事吗？\par

% structure: p0046 source=h3 target=subsection reason=relative-heading-rank; base=h2

\subsection{第十一章 基督徒妇女更无理由像外邦人那样盛装出现，相反，她们的外表应始终与她们有别}

此外，你们有何理由过分奢华地出现在公共场合呢？你们本已远离那些需要如此炫耀的场合。因为你们既不巡游神庙，也不要求出席公共表演，更与异教徒的节庆毫无关联。正是为了所有这些公共集会，以及频繁的看与被看，一切（服饰的）排场才展示于众目之下；要么是为了从事淫乐的买卖，要么是为了膨胀\Quote{虚荣}。\Emph{你们}却没有出现在公共场合的理由，除非是严肃的事。要么是探望患病的弟兄，要么是举行祭献，要么是宣讲天主的圣言。无论你提到哪一种，都是庄重圣洁的事务，不需要特别的服饰，既无需（刻意）修饰，也无需（放任）随便。如果外邦人的友谊和善意的职责要求你去，为何不穿戴你自己的铠甲出发？而且更应如此，因为你所面对的是不信主的人。这样，天主的婢女与魔鬼的婢女之间才有区别，使你能成为她们的榜样，她们也能因你而受到教诲；这样（如宗徒所说）\Quote{天主能在你们的身体上受显扬}。不过，天主是在\Emph{身体}上借着克己而受显扬；当然也借着适合克己的服饰受显扬。但有人会争辩说：\Quote{不要因我们使这名字受亵渎，如果我们从旧有的风格和服饰中做出任何贬损的改变。}那么，我们就不要摒弃旧日的恶习！如果我们必须保持同样的外表（如前），就让我们保持同样的品格吧；那样，外邦人就不会亵渎了！这真是伟大的亵渎，因为有人会说：\Quote{自从她成为基督徒，她就穿着更朴素的衣裳行走！}难道你会害怕从你变得更为富足之时起显得更贫穷，从你变得更为洁净之时起显得更污秽吗？基督徒的行走，是遵照外邦人的法令，还是遵照天主的法令呢？\par

% structure: p0048 source=h3 target=subsection reason=relative-heading-rank; base=h2

\subsection{第十二章 此类外在装饰实为娼妓之举，故不适于端庄妇女}

但愿我们不做正当亵渎的因由！然而，被称为端庄祭司的你，若是像\Emph{不}端庄的人那样涂脂抹粉、盛装出现在公共场合，那岂不是更能引起亵渎？否则，那些可怜的不幸的公共淫欲的受害者，还有什么不如你呢？尽管有些法律曾禁止她们使用婚姻和主妇的装饰，但如今这时代日益增长的腐化已使她们几乎与所有最尊贵的妇女平等，以至于难以区分。然而，圣经也提示我们：娼妓般的形体吸引力总是与身体的卖淫结合在一起，并适合于此。那统治七座山和众多河流的强大城邦，被主称为娼妓。但哪种装束是使她与这一称谓相比拟的工具呢？她当然\Quote{身穿紫色、朱红色，戴满黄金、宝石}。没有这些，一个受诅咒的娼妓如何能被描述！正是塔玛尔\Quote{涂脂抹粉，打扮自己}，才使犹大以为她是妓女；因此，因为她被\Quote{帕子}遮掩，她的装束品质让人误以为她是妓女——他便如此判断她，与她说话并讨价还价。由此我们进一步确认了教训：必须尽力防范一切不端的联想和猜疑。因为为何一个贞洁心灵的正直会被邻居的猜疑玷污？为何别人会指望我做出我所厌恶的事？为何我的装束不预先宣告我的品格，以避免我的精神因无耻通过耳朵而受伤？且允许我假设一个端庄妇女的外观：但假设一个\Emph{不}端庄妇女的外观，则至少是\Emph{不}许可的。\par

% structure: p0050 source=h3 target=subsection reason=relative-heading-rank; base=h2

\subsection{第十三章 仅天主知道我们贞洁还不够：我们在人前也应显得如此。尤其是在这迫害时期，我们应使身体习惯于可能遭受的艰辛}

也许有（妇女）会说：\Quote{我不必得到人的赞许；因为我并不需要人的证明：天主是内心的监察者。}（这一点）我们都知道；但条件是，我们记得同一（天主）曾通过宗徒所说的话：\Quote{\BibleQuote{让你们的正直显现在人前。}}为了什么目的，除非是使邪恶完全无法接近你们，或者你们可以成为恶人的榜样和见证？不然的话，（那句话）是什么：\Quote{\BibleQuote{让你们的善行发光？}}再者，为什么主称我们为世界的光；为什么祂把我们比作建在山上的城；如果我们不在黑暗中发光，并屹立于那些沉沦者之中呢？如果你把灯藏在斗底下，你必定会完全留在黑暗中，并且被许多人撞到。使我们成为世界之光的事物就是这些——我们的善行。而且，什么是\Emph{善}呢？只要是真实而丰盈的，就不爱黑暗：它乐于被看见，并因被人指指点点而欢欣。对于基督徒的谦逊，仅仅\Emph{是}谦逊是不够的，还要\Emph{显得}谦逊。因为它的丰盈应如此之大，以至于它可以从心灵流淌到服饰，从良心迸发到外表；这样，甚至从外面它也可以像凝视自己的装备一样——（一种装备）适于永远将信德作为其内在的居住者。因为那些因其柔软和柔弱而倾向于消除信德的刚毅的精致物品，应该被丢弃。否则，我不知道那习惯了被棕叶般的手镯环绕的手腕，是否能忍受变得如同自己的锁链般麻木僵硬！我不知道那因脚环而喜悦的腿，是否能忍受被挤入脚镣！我担心那被珍珠和翡翠的环套所围绕的脖子，是否会给大刀留出空间！因此，蒙福的（姊妹们），让我们默想艰难，这样我们就不会感到它们；让我们放弃奢华，这样我们就不会后悔它们。让我们准备好忍受一切暴力，没有什么我们害怕留下的东西。正是这些东西是阻碍我们望德的锁链。如果我们渴望天上的装饰，就让我们抛弃地上的装饰。不要爱金子；就是（这一种物质）中刻着以色列人民的一切罪过。你们应该\Emph{憎恨}那毁灭了你们祖先的东西；那被他们离弃天主时所敬拜的东西。甚至\Emph{那时}（我们发现）金子是火的食物。但基督徒始终，且现在比以往任何时候，他们的时光不是在金子中而是在铁中度过：殉道的长袍（现在）正在准备：那将要背负我们的天使（现在）正在被等候！你们出去（迎接她们）时，已经穿戴了先知和宗徒的化妆品和装饰；从纯朴中取得你们的洁白，从谦逊中取得你们的红润；用羞怯描绘你们的眼睛，用沉默描绘你们的嘴巴；将天主的话栽种在你们的耳朵里；将基督的轭套在你们的脖子上。向你们的丈夫顺服你们的头，这样你们就会被充分装饰。让你们的双手忙于纺织；让你们的双脚留在家里；这样你们会比（用金子装扮自己）更能\Quote{蒙悦纳}。用正直的丝绸、圣洁的细麻布、谦逊的紫布来装扮你们自己。这样装扮，你们将拥有天主作你们的爱人！\par

\SourceRecord{Translated by S. Thelwall.；Ante-Nicene Fathers；Vol. 4.；Edited by Alexander Roberts, James Donaldson, and A. Cleveland Coxe.；Buffalo, NY: Christian Literature Publishing Co.,；1885.；<http://www.newadvent.org/fathers/0402.htm>.}
